\documentclass[11pt]{article}

\usepackage[utf8]{inputenc}
\usepackage[T1]{fontenc}
\usepackage{lmodern}
\usepackage[margin=1in]{geometry}
\usepackage{amsmath,amssymb}
\usepackage{graphicx}
\usepackage{booktabs}
\usepackage{microtype}
\usepackage[round]{natbib}
\usepackage[hidelinks]{hyperref}

\bibliographystyle{plainnat}

% convenience macros
\newcommand{\R}{\mathbb{R}}
\newcommand{\Rnn}{\R_{\ge 0}}
\newcommand{\DIS}{D_{\mathrm{IS}}}
\newcommand{\DKL}{D_{\mathrm{KL}}}
\newcommand{\HSIC}{\mathrm{HSIC}}
\newcommand{\code}[1]{\texttt{#1}}

\title{Multiplicatively-coherent spatial pathway-activity analysis\\
of imaging mass spectrometry and spatial transcriptomics}

\author{%
  TBD\\
  \small SciLifeLab / KTH Royal Institute of Technology, Department of Gene Technology\\
  \small Correspondence: \texttt{lukas.kall@scilifelab.se}
}
\date{Working draft}

\begin{document}
\maketitle

\begin{abstract}
Imaging mass spectrometry (IMS / MALDI-MSI) yields, for each tissue pixel, a spectrum of
metabolite ion intensities, and a natural analysis goal is to decompose the image into a
small number of interpretable \emph{pathway-activity components} --- for each component, a
spatial map of \emph{where} a coordinated metabolic programme is active and a loading
vector of \emph{which} metabolites participate. The established approach factorises the
pixel\,$\times$\,$m/z$ matrix with PCA, ICA, or non-negative matrix factorisation (NMF),
usually after a log transform intended to stabilise the signal-dependent measurement
error. We argue this conventional pipeline contains an internal contradiction: a logarithm
can serve the \emph{error model} (turning multiplicative measurement noise into additive
noise that least squares handles) \emph{or} the \emph{coupling model} (preserving the
multiplicative co-regulation that defines a pathway), but not both. We present
\code{GAIT}, a framework for \emph{spatial pathway-activity analysis} built around
this distinction. It performs
Itakura--Saito (IS) or Kullback--Leibler (KL) NMF in \emph{linear} space, so each rank-1
component is literally a multiplicative outer product (coupling preserved) while a
scale-invariant divergence supplies the multiplicative-error model (no log needed);
detector saturation is handled by a masked loss that right-censors clipped entries. An
optional kernel (HSIC) penalty further lets the components be driven toward ICA-style
statistical (mutual-information) independence while remaining non-negative and
multiplicative. On synthetic data with known structure the engine recovers the planted
pathways and the noise-model diagnostic selects the correct divergence. On a real
catecholamine/serotonin MALDI-MSI section ($7{,}730$ pixels $\times$ $27$ metabolites) the
pipeline auto-selects the KL loss, decomposes the image into five spatially-distinct
components explaining ${\sim}85\%$ of the signal, and---without prompting---surfaces three
pairs of metabolites that are perfectly co-linear because they are unresolved isobars at
$5$~ppm, illustrating why downstream pathway labelling must be ambiguity-aware. The same
linear-space formulation transfers beyond mass spectrometry: on a mouse-brain
spatial-transcriptomics section a single neurotransmission pathway decomposes into
anatomically distinct \emph{dominating} components---striatal versus cortical---recovering
brain regions from gene counts without a log transform. The
contribution is not a new factorisation algorithm but the \emph{statistical coherence of
the pipeline}: matching a multiplicative coupling model to a multiplicative error model,
with saturation censoring and a path to isomer-aware pathway enrichment.
\end{abstract}

\section{Introduction}
\label{sec:intro}

Spatial biology profiles, at many locations across intact tissue, the abundances of many
molecular species at once. Imaging mass spectrometry (IMS / MALDI-MSI) records hundreds of
metabolite or lipid ions per pixel; spatial transcriptomics records thousands of gene counts
per spot; spatial proteomics and related assays add further feature types on the same
spatial substrate. Across these technologies the data share one structure --- a non-negative
matrix of (location $\times$ feature) abundances registered to tissue coordinates --- and one
recurring goal: to summarise a section as a handful of interpretable \emph{spatial
pathway-activity components}, each pairing a spatial map of \emph{where} a coordinated
molecular programme is active with a loading vector of \emph{which} features participate.
Nothing in this goal, nor in the method we develop for it, is tied to a particular
instrument: we validate it on two modalities --- IMS metabolites and spatial-transcriptomic
counts --- but it applies to any spatial-omics matrix whose features are non-negative,
multiplicatively co-regulated, and measured with signal-dependent error.

Matrix factorisation of the (locations\,$\times$\,features) matrix is the established
multivariate tool, and it is most thoroughly explored in imaging mass spectrometry. PCA,
ICA, and NMF have all been applied to MALDI-MSI,
and the comparison by \citet{siy2007} on mouse cerebellum found ICA and NMF more effective
than PCA, with NMF's non-negativity matching the physical non-negativity of ion abundances
and yielding parts-based, additive factors. \citet{leuschner2019} make the
spectra-plus-image reading explicit: one factor holds mass spectra, the other their spatial
distribution as pseudo-channels, and non-negativity supports biological interpretability
--- exactly the ``pathway activity graph $+$ pathway activity image'' pairing we adopt.
None of these components \emph{is} a pathway; the pathway interpretation is imposed
afterwards by annotation and enrichment \citep{jones2014,wittmann2025}.

Yet the dominant pathway-analysis paradigm asks a narrower question than the spatial setting
invites. Over-representation analysis (ORA) tests whether a pathway's members are unexpectedly
frequent among the features called differential, and gene-set enrichment analysis (GSEA)
tests whether they are shifted coherently towards one end of a ranked list
\citep{subramanian2005}; what null hypothesis such tests actually address --- and whether it
is the biologically intended one --- is itself a subtle, much-discussed question
\citep{goeman2007}. Both are well validated for detecting a \emph{quantitative, overall}
change in a pathway's regulation between conditions, and a decade of benchmarking has refined
them for exactly that \citep{khatri2012}. But each reduces a pathway to a single enrichment
statistic, and so cannot express \emph{how} a pathway is regulated --- which members carry the
activity, and whether distinct sub-programmes of one pathway are active in distinct places. In
a spatial context this is usually the question of interest: not whether a pathway is globally
up or down, but how its activity is \emph{organised across the tissue}. Our mouse-brain
example (Section~\ref{sec:spatial}) is precisely this case --- a single neurotransmission
pathway whose \emph{dominating} sub-programme switches between striatum and cortex, a
distinction no per-pathway enrichment score can represent.

We therefore pose \emph{spatial pathway-activity analysis}: rather than score a pathway, we
\emph{decompose} it into spatial components, each a \emph{where} (a tissue activity map)
paired with a \emph{which} (a feature-loading vector), so that region-specific sub-programmes
become first-class objects. Estimating a sample-level pathway \emph{activity} --- instead of
testing its genes for enrichment --- has precedent: such activities are stronger prognostic
markers than the individual transcripts \citep{jeuken2022}, and their association with a
phenotype can be scored by mutual information without gene-level comparisons
\citep{jeuken2024}; GAIT carries this paradigm into tissue space and makes it
multiplicatively coherent. What such a decomposition should optimise is fixed by two
properties of the data that one conventional choice --- the log transform --- fatally
entangles. \textbf{(R2)} Pathway coupling is \emph{multiplicative}: co-regulated features
scale together, so a pathway is a rank-one outer product $u\,v^{\!\top}$ of a per-location
activity and a per-feature loading. \textbf{(R1)} Measurement error is \emph{multiplicative}
as well --- signal-dependent gain noise in mass spectrometry, sampling noise in sequencing
--- which the field tames by taking logarithms, making the error additive so least squares
applies. But the logarithm cannot serve both at once: because $\log(a\,s) = \log a + \log s$,
it converts the shared multiplicative activity of a pathway into a mere \emph{additive}
offset, so a linear factorisation of logged data no longer studies multiplicative coupling
(R2) --- even though the log was taken only to fix the error (R1). \code{GAIT} resolves
the conflict instead of splitting it: it keeps the data \emph{linear}, so the outer product
remains genuine multiplicative coupling, and recovers the multiplicative-error property from a
\emph{scale-invariant loss} (Itakura--Saito or Kullback--Leibler) rather than from a transform
--- satisfying R1 and R2 simultaneously. We give the full argument in the Supplementary
Material (\code{supplement.tex}); Section~\ref{sec:methods} develops the resulting method, a
linear-space IS/KL non-negative factorisation with optional saturation censoring and an
optional ICA-style independence penalty.

\section{Methods}
\label{sec:methods}

\subsection{Data model and notation}
\label{sec:notation}

We arrange the data as $X \in \Rnn^{P\times M}$ with $P$ pixels (rows) and $M$ metabolite
ions (columns), and seek a rank-$K$ factorisation
\begin{equation}
  X \approx U V, \qquad U \in \Rnn^{P\times K}\ \text{(spatial scores)}, \qquad
  V \in \Rnn^{K\times M}\ \text{(spectral loadings)}.
\end{equation}
Each component $k$ is the pair $(U_{:,k},\,V_{k,:})$: $V_{k,:}$ is the \textbf{pathway
activity graph} (a loading over the metabolite network) and $U_{:,k}$ is the \textbf{pathway
activity image} (a per-pixel activity map placed back on the tissue grid). The rank-1 term
$U_{:,k} \otimes V_{k,:}$ is interpreted as one pathway.

\subsection{Itakura--Saito / KL NMF in linear space}
\label{sec:route2}

We minimise an elementwise divergence between $X$ and the reconstruction $Y = UV$ subject to
$U, V \ge 0$. The Itakura--Saito divergence,
\begin{equation}
  \DIS(x, y) = \frac{x}{y} - \log\frac{x}{y} - 1,
\end{equation}
is \textbf{scale-invariant}, $\DIS(\lambda x, \lambda y) = \DIS(x, y)$, so it penalises the
\emph{ratio} $x/y$ and a factor-of-two reconstruction error costs the same at high and low
abundance. This is the property the log transform was meant to deliver, obtained here
without leaving linear space, so R1 and R2 hold simultaneously. IS corresponds to
multiplicative Gamma-type gain noise (variance $\propto$ mean$^2$). When the dominant noise
is instead ion-counting/shot noise (variance $\propto$ mean) the matched choice is the
generalised Kullback--Leibler divergence,
\begin{equation}
  \DKL(x, y) = x\,\log\frac{x}{y} - x + y,
\end{equation}
which is Poisson in spirit. Both are exposed behind \code{loss="is"} and \code{loss="kl"};
the choice is made empirically (Section~\ref{sec:diagnostics}).

\paragraph{Saturation as right-censoring.} High-abundance ions can saturate the detector;
the clipped intensity then re-emerges in later components as a spurious ``compensation''
factor. We treat saturated $(\text{pixel}, \text{ion})$ entries as \emph{missing} via a
binary weight matrix $W \in \{0,1\}^{P\times M}$ and minimise the masked objective
\begin{equation}
  \min_{U, V \ge 0}\ \sum_{p,m} W_{pm}\, D\!\big(X_{pm},\, (UV)_{pm}\big).
\end{equation}
Because IS and KL are defined directly on non-negative linear data, the mask composes
cleanly --- a practical advantage of the linear route over a log$+$MSE pipeline, which would
first need an intermediate $\operatorname{asinh}$ transform under which neither the additive
nor the multiplicative argument holds exactly. The mask is built from per-ion intensity
histograms: a pile-up spike at a ceiling indicates hard clipping and is masked; a smoothly
bending tail indicates soft compression and masking is optional. Masking is \textbf{off by
default} and the engine warns when a pile-up is detected.

\paragraph{Optimisation.} We use multiplicative-update (MU) rules --- Lee--Seung form for
Frobenius/KL \citep{lee2001}, F\'evotte--Bertin--Durrieu form for IS \citep{fevotte2009} ---
with the weight matrix $W$ carried inside both the numerator and the denominator of every
update so that censored entries influence neither the fit nor its normalisation. IS-NMF is
non-convex and MU can be unstable, so IS fits are \textbf{warm-started from a few KL
iterations}, every quantity is floored by a small $\varepsilon$ to avoid division by zero,
and the best of several random restarts (\code{n\_init}) is kept. Outputs $U$ and $V$ are on
interpretable linear-abundance units.

\subsection{Illustration layer}
\label{sec:illustration}

Each fitted component is rendered as the two objects above. The pathway activity image
reshapes $U_{:,k}$ onto the acquisition grid as a heatmap, supporting non-rectangular tissue
through an explicit pixel-index map. The pathway activity graph colours a metabolite-reaction
network; colouring nodes by the raw loading $V_{k,:}$ is uninformative because metabolite
concentrations span several orders of magnitude, so a handful of abundant ions saturate the
colour scale and the rest appear uniformly dark. We instead colour each metabolite by the
\emph{fraction of its own variation that the component explains},
\begin{equation}
  F_{k,m} \;=\; \frac{V_{k,m}^2\,\operatorname{Var}_p\!\big(U_{:,k}\big)}
                     {\sum_{k'} V_{k',m}^2\,\operatorname{Var}_p\!\big(U_{:,k'}\big)}
            \;\in[0,1],
  \label{eq:explained}
\end{equation}
the variance of component $k$'s rank-1 contribution $U_{:,k}V_{k,m}$ to metabolite $m$,
normalised per metabolite so that the columns sum to one over components. The product
$V_{k,m}\lVert U_{:,k}\rVert$ is invariant to the NMF scale ambiguity, so $F_{k,m}$ is well
defined; it places every metabolite on a common $[0,1]$ scale, so a low-abundance ion
``owned'' by a component reads as brightly as an abundant one. (As the components are not
orthogonal, $F$ is a normalised share rather than an exact variance partition, exact in the
independent-component limit of Section~\ref{sec:independence}.)

The same construction has a spatial dual, obtained by swapping the roles of pixels and
metabolites: the per-pixel fraction of variation that component $k$ explains is
\begin{equation}
  G_{p,k} \;=\; \frac{U_{p,k}^2\,\operatorname{Var}_m\!\big(V_{k,:}\big)}
                     {\sum_{k'} U_{p,k'}^2\,\operatorname{Var}_m\!\big(V_{k',:}\big)}
            \;\in[0,1],
  \label{eq:spatial-explained}
\end{equation}
with each pixel's shares summing to one over components.

\paragraph{Two complementary spatial readings.} The raw activity image $U_{:,k}$ and the
spatial fraction-explained image $G_{:,k}$ answer \emph{different} questions and are both
worth showing. $U_{:,k}$ is the component's absolute activity --- \emph{where the pathway is
active}; $G_{:,k}$ is the component's \emph{local share} of the signal --- \emph{where it
dominates the variation}, after discounting pixels where a more active component also fires.
They are not interchangeable: a component that is broadly but weakly active has a diffuse
$U_{:,k}$ yet a sharp $G_{:,k}$ confined to the region it actually controls, while a region
of high absolute activity that is co-explained by several components is bright in $U$ but
only moderate in $G$. We therefore present both, alongside the graph fraction-explained
$F_{k,:}$ (the per-metabolite analogue of $G$). A single bridge call
(\code{illustrate\_component}) produces the graph and image views from a fitted result, so
the decomposition and visualisation layers are decoupled but composable.

\subsection{Diagnostics}
\label{sec:diagnostics}

Four checks precede interpretation. (i)~A \textbf{variance-vs-mean} fit of $\log\mathrm{Var}$
against $\log\mathrm{Mean}$ across ions estimates the noise exponent and recommends the loss:
slope $\approx 2 \to$ IS, $\approx 1 \to$ KL, $\approx 0 \to$ Frobenius. (ii)~A
\textbf{saturation/mechanism} check inspects per-ion histograms for clipping pile-ups.
(iii)~A \textbf{compensation-artifact} check verifies that a suspected saturating ion's
reappearance in later components is spatially anti-correlated with its primary-component
image (i.e.\ a measurement nonlinearity, not co-localisation). (iv)~A
\textbf{component-recovery} check confirms that masking/IS collapses spurious compensating
components relative to a Frobenius baseline.

\subsection{ICA-style independent components: a mutual-information penalty}
\label{sec:independence}

Non-negative factorisation --- including the IS/KL variant above --- yields parts-based but
generally \emph{correlated} components; it does not impose the statistical
\textbf{independence} that independent component analysis (ICA) targets, and which is often
desirable for higher components (``after the dominant programme, what is the next
\emph{independent} one?''). Classical ICA achieves independence by whitening and rotating in
signed space, which would destroy both the non-negativity and the multiplicative
outer-product structure central to our principles. We therefore keep the IS/KL fidelity term
and \emph{add} a penalty that targets mutual information directly:
\begin{equation}
  \min_{U, V \ge 0}\ D(X \,\|\, UV)\ +\ \lambda \sum_{i<j} \HSIC\!\big(U_{:,i},\, U_{:,j}\big).
\end{equation}
The Hilbert--Schmidt Independence Criterion (HSIC) \citep{gretton2005} with a characteristic
(Gaussian RBF) kernel is zero \emph{if and only if} the two variables are independent ---
equivalently, iff their mutual information vanishes --- so minimising the pairwise HSIC
drives the components toward MI-sense orthogonality. This is precisely the kernel dependence
contrast used by Kernel ICA \citep{bach2002}; the novelty here is only that we apply it as a
\emph{penalty on a non-negative factor} rather than as an orthogonal-rotation objective, so
independence is obtained without leaving the multiplicative, non-negative model. A linear
kernel reduces the penalty to ordinary second-order decorrelation, which is \textbf{not}
sufficient for independence; the RBF kernel captures dependence of all orders.

To scale to many pixels we approximate the RBF kernel and its gradient with random Fourier
features \citep{rahimi2007}, giving the penalty in $O(P\,D)$ for $D$ features. Optimisation
warm-starts from the plain solution, then alternates a masked multiplicative-update step for
$V$ with a projected-gradient $+$ backtracking step for $U$ on the deterministic
$(\text{data} + \lambda\cdot\text{penalty})$ objective, projecting onto non-negativity.
Because the raw HSIC is many orders of magnitude smaller than the divergence, $\lambda$ is
made dimensionless by rescaling the penalty gradient to the data gradient's norm at each
step: $\lambda = 0$ recovers the plain fit and $\lambda \approx 1$ pushes independence about
as hard as the data fidelity. Independence is penalised on the spatial maps $U$ (the
spatial-ICA analogue) by default, or on the loadings $V$. We verify the outcome with an
independent diagnostic --- the normalised HSIC (centred kernel alignment, in $[0, 1]$)
between component pairs, computed with a direct RBF kernel.

\section{Results}
\label{sec:results}

\subsection{Synthetic validation}
\label{sec:synthetic}

We generated a $24 \times 24$-pixel synthetic image with two spatial activity regions
(Gaussian blobs), each multiplicatively coupling a disjoint set of six metabolites, plus
weak background ions, corrupted with multiplicative Gamma noise and one hard-clipped
(saturating) ion. The variance-vs-mean diagnostic recovered a near-quadratic noise law and
recommended IS; the saturation detector flagged the clipped ion and built the corresponding
mask. Masked IS-NMF recovered both planted components: the spatial maps matched the two
blobs and the top-loading ions of each component were exactly the planted members. A
unit-tested component-recovery check confirms that masking does not increase the
off-component spread of the saturating ion relative to an unmasked baseline.

\subsection{Log vs.\ linear: recovering the multiplicative coupling}
\label{sec:benchmark}

To test the central claim directly --- that a linear-space factorisation with a
multiplicative-error loss recovers a pathway's multiplicative structure where the
conventional log pipeline does not --- we planted $K = 3$ region-structured components, each
multiplicatively coupling a group of features whose loadings span three orders of magnitude
(so the log's dynamic-range compression is exercised), corrupted them with multiplicative
Gamma noise, and asked five methods to recover the planted structure
(Table~\ref{tab:benchmark}). Recovery of the spatial maps ($U$) and of the loadings ($V$) is
scored as the mean absolute correlation between each recovered factor and its
optimally-matched planted factor; \emph{region accuracy} is the accuracy of the per-location
dominant-component map against the planted regions.

The result isolates the two design choices. \textbf{Coupling fidelity (R2)} is read from
$V$-recovery: every \emph{linear} method recovers the multiplicative loadings essentially
perfectly ($\approx 1.00$), whereas $\log(1{+}x)$ before NMF drops to $0.79$ and before PCA to
$0.54$. Tellingly, $\log(1{+}x){+}$NMF still recovers \emph{where} each component is active
($U$-recovery $0.90$, region accuracy $0.96$) but not \emph{which} features couple or
\emph{how strongly} ($V$-recovery $0.79$): the log leaves the spatial activity as a
recoverable additive offset while degrading exactly the multiplicative loading structure that
defines the pathway --- the precise failure anticipated in Section~\ref{sec:intro}.
\textbf{Error-model matching (R1)} is read from the top rows: under multiplicative noise the
Itakura--Saito loss is best (region accuracy $0.98$), and when the same experiment is re-run
with Poisson (count) noise the ranking flips and KL leads ($U/V/\text{region} = 1.00/1.00/0.98$,
IS falling to $0.84$) --- so the variance-vs-mean diagnostic (Section~\ref{sec:diagnostics})
selects the better loss in each regime. The conventional $\log{+}$PCA baseline is worst on
every comparable metric (for the signed, non-parts-based PCA factors a dominant-component map
is ill-defined, so its region accuracy is not directly comparable; the $U/V$ correlations,
$0.64/0.54$, are). The benchmark is reproducible with \code{python -m gait.benchmark}.

\begin{table}[t]
  \centering
  \caption{Recovery of planted multiplicative, region-structured components under
  multiplicative (Gamma) noise ($K=3$): mean $|\text{corr}|$ of each recovered factor to its
  matched ground-truth factor (for $U$ and $V$), and dominant-component accuracy vs the
  planted regions (higher is better). Linear methods recover the loadings $V$ essentially
  perfectly; both log pipelines do not. \code{python -m gait.benchmark} regenerates this and
  the Poisson-noise variant discussed in the text.}
  \label{tab:benchmark}
  \begin{tabular}{lccc}
    \toprule
    method & $U$-recovery & $V$-recovery & region acc. \\
    \midrule
    GAIT (IS)              & 0.99 & \textbf{1.00} & \textbf{0.98} \\
    GAIT (KL)              & 0.96 & \textbf{1.00} & 0.96 \\
    linear NMF (Frobenius) & 0.94 & \textbf{1.00} & 0.95 \\
    $\log(1{+}x)$ + NMF    & 0.90 & 0.79 & 0.96 \\
    $\log(1{+}x)$ + PCA    & 0.64 & 0.54 & 0.02 \\
    \bottomrule
  \end{tabular}
\end{table}

\subsection{Real MALDI-MSI section (catecholamine/serotonin metabolites)}
\label{sec:realdata}

We applied the pipeline to a single tissue section (\code{01\_pd\_51}) profiled at $5$~ppm,
comprising \textbf{$7{,}730$ pixels $\times$ $27$ metabolites} of the catecholamine and
serotonin pathways (dopamine, L-DOPA, HVA, serotonin, 5-HIAA, \ldots) with pixel coordinates
spanning a $62 \times 144$ grid (the tissue is non-rectangular; $7{,}730$ of $8{,}928$ grid
cells are occupied). The matrix is non-negative with no missing values and is $45.2\%$ exact
zeros, with a dynamic range from $0$ to ${\sim}1.8 \times 10^6$ --- consistent with strongly
signal-dependent error.

\paragraph{Automatic model choices.} The variance-vs-mean diagnostic returned a noise
exponent of \textbf{$1.35$}, recommending the \textbf{KL} loss (intermediate between Poisson
and fully multiplicative, and the stable choice given the large fraction of zeros, which IS
penalises harshly). No histogram pile-up was detected on any ion, so no entries were masked.
The elbow heuristic over $K \in \{2,\ldots,8\}$ selected \textbf{$K = 5$}.

\paragraph{Decomposition quality.} KL-NMF at $K = 5$ (six restarts) reconstructed the image
with a relative Frobenius error of \textbf{$0.154$}, i.e.\ roughly \textbf{$85\%$ of the
signal energy} captured by five components. The components are spatially distinct, occupying
different bands of the section rather than re-describing one dominant region. This is
clearest in the \textbf{first three components} (Figure~\ref{fig:overview}): component~0 is
serotonin-pathway--dominated (Serotonin, 5-HIAA) together with the high-abundance terminal
catabolites HVA/MOPEGAL and 3-OMD, and localises to the lower portion of the section;
component~1 (DOPEG, 5-HIAA, HVA/MOPEGAL) is comparatively diffuse; and component~2 (3-OMD,
MOPEG, 3-MT) concentrates in a distinct upper band, with 3-OMD by far its strongest loading
(Figure~\ref{fig:overview}, right column). Their dominant metabolites
(Table~\ref{tab:loadings}) are chemically coherent, and the later components continue the
pattern: component~3 is a clean amine-catabolite group (DOPAC/DOPEGAL,
Epinephrine/Normetanephrine, Metanephrine) and component~4 is precursor-dominated
(L-DOPA, 3-MT).

The two spatial readings of Section~\ref{sec:illustration} are complementary here rather than
redundant (Figure~\ref{fig:overview}, top two rows). Component~1 is the clearest case: its raw
activity $U_{:,1}$ is broad and low-contrast across much of the section, but its
fraction-explained map $G_{:,1}$ sharpens to the mid-section territory it actually dominates,
once the lower band (owned by component~0) and the upper band (component~2) are discounted.
Conversely, the high-activity lower band of component~0 remains the region it both occupies
and dominates, so $U_{:,0}$ and $G_{:,0}$ largely agree. Presenting both therefore separates
``where a pathway is active'' from ``where it accounts for the local signal,'' which a single
map conflates.

\begin{figure}[tbp]
  \centering
  \includegraphics[width=\textwidth]{figures/components_overview.png}
  \caption{\textbf{First three components of \code{01\_pd\_51}: three complementary views per
  component.} \emph{Row 1 --- activity score:} the per-pixel score $U_{:,k}$ for components
  0--2, placed back on the $62 \times 144$ tissue grid (KL-NMF, $K = 5$) and titled with the
  component's top-3 metabolites; this shows \emph{where the pathway is active}. \emph{Row 2
  --- spatial fraction explained:} the per-pixel share $G_{:,k}\in[0,1]$
  (Equation~\ref{eq:spatial-explained}); this shows \emph{where the component dominates the
  local signal}, sharpening diffuse activity (clearest for component~1) to the territory it
  actually controls. \emph{Row 3 --- graph fraction explained:} the pathway activity graph
  over the catecholamine/serotonin reaction network, with each metabolite coloured by the
  fraction $F_{k,m}\in[0,1]$ of its variation explained by the component
  (Equation~\ref{eq:explained}) rather than by the raw loading, which removes the
  concentration imbalance that otherwise renders the colouring near-binary (hand-laid layout,
  labels haloed for contrast). Rows~1 and~2 are deliberately both shown because they answer
  different questions (Section~\ref{sec:illustration}). Regenerate from the source parquet
  with \code{python manuscript/make\_figures.py <parquet>}.}
  \label{fig:overview}
\end{figure}

\paragraph{An ambiguity surfaced by the data.} Three metabolite pairs ---
\textbf{DOPAC/DOPEGAL}, \textbf{HVA/MOPEGAL}, and \textbf{Epinephrine/Normetanephrine} ---
have \emph{perfectly identical} intensity columns and therefore always co-load on the same
component with equal weight. These are unresolved isobars/isomers at MS1 $5$~ppm: the
factorisation correctly reports that they are indistinguishable from these data alone. This
is a concrete, unprompted instance of the MS1 (MSI Level~2) ambiguity that makes na\"ive
molecule-name pathway labelling invalid, and it motivates ambiguity-aware enrichment
(Section~\ref{sec:discussion}) rather than a one-molecule-per-ion assignment.

\begin{table}[t]
  \centering
  \caption{Top-5 metabolites by loading for each KL-NMF component ($K = 5$,
  \code{01\_pd\_51}). HVA$\equiv$MOPEGAL, DOPAC$\equiv$DOPEGAL, and
  Epinephrine$\equiv$Normetanephrine are unresolved isobar pairs; their equal loadings
  reflect identical measured columns.}
  \label{tab:loadings}
  \begin{tabular}{cl}
    \toprule
    Comp & Top-loading metabolites \\
    \midrule
    0 & Serotonin, HVA, MOPEGAL, 5-HIAA, 3-OMD \\
    1 & DOPEG, 5-HIAA, HVA, MOPEGAL, Serotonin \\
    2 & 3-OMD, MOPEG, 3-MT, 5-HTOL, DOPAL \\
    3 & DOPEGAL, DOPAC, Epinephrine, Normetanephrine, Metanephrine \\
    4 & L-DOPA, 3-MT, DOMA, DOPEG, DOPAL \\
    \bottomrule
  \end{tabular}
\end{table}

\subsection{ICA-style independence}
\label{sec:ica}

We tested the independence penalty (Section~\ref{sec:independence}) in two regimes. On a
synthetic problem with two strongly but \emph{non-linearly} dependent non-negative sources
(so they are nearly uncorrelated yet far from independent), plain KL-NMF recovered components
with mean off-diagonal normalised HSIC of $0.18$. Adding the penalty drove this toward zero
monotonically with $\lambda$, at a small and controllable reconstruction cost
(Table~\ref{tab:independence}): at $\lambda = 1$ dependence fell to $0.08$ for essentially no
loss of fit, and by $\lambda = 3$ the components were independent (nHSIC $= 0.00$) at a
relative-reconstruction cost rising from $0.146$ to $0.174$. This is the ICA trade-off ---
exchanging a little reconstruction for statistical independence --- realised inside the
non-negative, multiplicative model.

On the real \code{01\_pd\_51} section the KL-NMF components were \textbf{already
near-independent} (mean off-diagonal normalised HSIC $\approx 0.05$), so the penalty had
little to do: dependence was essentially unchanged and, reassuringly, the reconstruction
error was preserved ($0.154 \to 0.155$) even at $\lambda = 8$. The penalty thus acts as a
safeguard/diagnostic when the data already separate into independent spatial programmes, and
as an active tool when they do not.

\begin{table}[t]
  \centering
  \caption{Independence penalty on two non-linearly dependent synthetic sources ($K = 2$, KL
  loss). $\lambda = 0$ reproduces the plain fit.}
  \label{tab:independence}
  \begin{tabular}{ccc}
    \toprule
    $\lambda$ & mean off-diag normalised HSIC & relative reconstruction error \\
    \midrule
    0 (plain KL-NMF) & 0.18 & 0.146 \\
    1  & 0.08 & 0.148 \\
    3  & 0.00 & 0.174 \\
    10 & 0.00 & 0.178 \\
    \bottomrule
  \end{tabular}
\end{table}

\subsection{Generality: region-specific pathway components in spatial transcriptomics}
\label{sec:spatial}

The error-vs-coupling argument is not specific to mass spectrometry. In spatial
transcriptomics the rows are tissue spots and the columns are gene counts; counting/sampling
noise is multiplicative in the sense of R1 (variance growing with the mean, Poisson to
over-dispersed) and pathway co-regulation is multiplicative in the sense of R2, so the
ubiquitous ``$\log(1+x)$-normalise, then PCA/NMF'' recipe commits exactly the conflation of
Section~\ref{sec:intro}. We therefore treat counts like the metabolites: a \emph{linear}
per-spot library-size normalisation (no log), then KL/IS-NMF in linear space.

We applied this to a mouse-brain Visium section (GEO \texttt{GSE232910}, sample
\code{V11L12-038\_A1}; $2{,}856$ spots $\times$ $32{,}285$ genes, raw counts) restricted to a
$36$-gene neuroactive-ligand/receptor (\emph{neurotransmission}) pathway, posing a question
the single-region IMS sections cannot: does a \emph{single} pathway carry \emph{different
dominating sub-programmes in different tissue regions}? KL-NMF at $K = 5$ resolves the
pathway into spatially segregated modules, and their \textbf{dominant-component map}---the
per-spot argmax of the fraction-explained $G_{p,k}$ (Equation~\ref{eq:spatial-explained})
---reconstructs the anatomy (Figure~\ref{fig:visium}). Component~0 is a striatal
medium-spiny-neuron programme
(\textit{Gpr88}, \textit{Ppp1r1b}/DARPP-32, \textit{Penk}, \textit{Adcy5}, \textit{Drd1});
component~1 is cortical-glutamatergic (\textit{Slc17a7}, \textit{Camk2a}, \textit{Satb2},
\textit{Grin2a}); a GABAergic module (\textit{Gad1}/\textit{Gad2}/\textit{Slc32a1}) and a
thalamic one (\textit{Tcf7l2}) follow. Against the dataset's own anatomical annotation,
component~0 dominates $71\%$ of striatum spots versus $9\%$ elsewhere, and component~1
dominates $41\%$ of non-striatum spots versus $1\%$ of striatum. Restricting the pathway to
the two dopamine receptors \textit{Drd1}/\textit{Drd2} sharpens the striatal localisation
(striatal vs.\ non-striatal activity ratio $3.5$; Mann--Whitney $p \approx 10^{-162}$ against
the region labels). The cell-body/terminal dissociation expected of dopaminergic biology is
also recovered: the synthesis enzymes \textit{Th}/\textit{Ddc}, whose cell bodies lie in the
midbrain outside this coronal plane, show \emph{no} striatal enrichment, whereas the
striatal receptor/terminal markers do.

Crucially, decomposing the \emph{same} data after the conventional $\log(1+x)$ transform
leaves the regional split intact but \emph{inflates} it---component~0's striatal dominance
reads $84\%$ under the log against $71\%$ under the linear KL fit---so the log preprocessing
exaggerates the separation while the structure itself is robust to the choice. The finding is
thus a property of the pathway, not of the preprocessing, and the linear KL fit (consistent
with the metabolite analysis) is the more faithful model. The analysis is reproduced by
\code{demo\_visium\_dopamine.py} on any Space~Ranger / GEO sample.

\begin{figure}[tbp]
  \centering
  \IfFileExists{figures/visium_dopamine_components.png}{%
    \includegraphics[width=\textwidth]{figures/visium_dopamine_components.png}%
  }{%
    \fbox{\parbox[c][3.2cm][c]{0.9\textwidth}{\centering\itshape
      \texttt{figures/visium\_dopamine\_components.png} not found --- regenerate with
      \texttt{python manuscript/make\_figures.py visium MATRIX.h5
      TISSUE\_POSITIONS.csv REGION.csv}.}}%
  }
  \caption{\textbf{A single neurotransmission pathway decomposes into region-specific
  dominating components on a mouse-brain Visium section} (\code{V11L12-038\_A1}, GEO
  \texttt{GSE232910}; raw counts, library-size normalised in linear space, KL-NMF $K=5$).
  \emph{Left ($K$ panels):} the per-spot activity score $U_{:,k}$ for each component, titled
  by its top-loading gene --- component~0 (\textit{Gpr88}/\textit{Ppp1r1b}/\textit{Penk})
  fills the striatum, component~1 (\textit{Camk2a}/\textit{Slc17a7}) the cortex, with
  GABAergic (\textit{Gad1}) and other modules following. \emph{Centre:} the
  \textbf{dominant-component} map, the per-spot argmax of the fraction-explained $G_{p,k}$
  (Equation~\ref{eq:spatial-explained}); a single pathway thereby segments the section into
  anatomically coherent territories. \emph{Right:} the dataset's striatum annotation, which
  the component-0 territory recovers ($71\%$ of striatum spots vs.\ $9\%$ elsewhere).
  Regenerate with \code{python manuscript/make\_figures.py visium <sample files>}.}
  \label{fig:visium}
\end{figure}

%% ---- PROVISIONAL (added ahead of the matched FMP-10 metabolite data; may be revised or
%%      removed once that analysis lands). Numbers from the GSE232910 cohort scan. ----
\subsection{Reproducibility, specificity, and an unbiased pathway scan}
\label{sec:robustness}

We asked whether the region-specific decomposition is reproducible and specific across the
mouse-brain cohort (GEO \texttt{GSE232910}; 19 Visium sections). Each section is summarised by
the best striatum-vs-rest AUC over the $K=5$ component activities (the most striatum-enriched
component). Across the six striatal sections the dopaminergic and the broad neurotransmission
programmes localise to the striatum reproducibly (median AUC $0.90$ and $0.95$;
Table~\ref{tab:robustness}); the GABAergic programme moderately ($0.72$ --- the striatal
projection neurons are themselves GABAergic); while the serotonergic and noradrenergic
programmes, whose neurons are not striatal, score at or near chance ($0.62$, $0.53$). This
\emph{specificity} --- rather than a uniformly high score --- is the evidence that GAIT
recovers genuine regional biology and not an artefact of the decomposition.

Scored instead against the intact-vs-lesioned annotation of the 6-OHDA model, the transcript
components sit at chance (AUC $\approx 0.55$ for the dopaminergic, MSN and neuropeptide sets).
This is expected, and motivating: the lesion depletes dopamine --- the \emph{metabolite} ---
in the lesioned striatum, but leaves the post-synaptic receptor \emph{transcripts}
(\textit{Drd1}/\textit{Drd2}) largely intact, so the transcriptomic programme cannot see the
lesion. It yields a concrete prediction: the matched dopamine-metabolite component (MSI) should
separate the lesioned hemisphere where the transcript component does not --- a
transcript-versus-metabolite dissociation only a multimodal measurement can resolve, which we
will test on the matched FMP-10 data.

Because the score underlying these maps requires no region annotation --- a
dominant-component-weighted Moran's I of the component activities times the entropy of the
dominant-component map --- it can rank \emph{arbitrary} pathways by how regionally structured
their activity is. We scanned the 304 KEGG pathways (\code{scan\_pathways.py}) across the eight
striatal sections of the cohort and aggregated each pathway by its \emph{median} score and
median rank over the sections. The top of the ranking is consistently occupied by the
\emph{synaptic / neuro-signalling} pathways --- \textbf{glutamatergic and dopaminergic synapse,
cAMP and calcium signalling, retrograde-endocannabinoid signalling} and neuroactive
ligand--receptor interaction (each present and scored in all eight sections) --- so the screen
recovers the striatal neurotransmission biology with no annotation and reproduces the supervised
ranking above, two independent metrics in agreement. The \emph{class} of top pathways is stable
section-to-section even where the exact order is not: the GABAergic-synapse set, top-ranked on
the single section first examined, falls to a median rank of $12$ across the eight, whereas the
dopaminergic synapse holds a median rank of $4$. (KEGG gene-set overlap means several lower
hits, e.g.\ the addiction pathways, largely re-detect the same synaptic programmes rather than
independent biology.)

\begin{table}[t]
  \centering
  \caption{Reproducibility and specificity across the six striatal sections of
  \texttt{GSE232910} (median best striatum-vs-rest AUC over $K=5$ GAIT components; linear
  counts, KL-NMF). Striatal systems localise reproducibly; non-striatal systems score at
  chance.}
  \label{tab:robustness}
  \begin{tabular}{lc}
    \toprule
    pathway & striatum AUC (median) \\
    \midrule
    neurotransmission (broad) & 0.95 \\
    dopaminergic              & 0.90 \\
    GABAergic                 & 0.72 \\
    serotonergic              & 0.62 \\
    noradrenergic             & 0.53 (chance) \\
    \bottomrule
  \end{tabular}
\end{table}
%% ---- end PROVISIONAL ----

\section{Discussion}
\label{sec:discussion}

\paragraph{Relation to prior work.} Plain matrix factorisation treats pixels as exchangeable
and discards spatial adjacency; \citet{fernsel2021} adds spatial coherence through a
total-variation penalty, but as an orthogonal-NMF \emph{clustering} method scored on
segmentation metrics, not a pathway recovery method --- a graph/TV penalty on the spatial
factor is a natural optional extension to the factorisation here, not part of the core. The
``component $\to$ pathway'' step is classically handled by annotate-then-enrich:
\citet{jones2014} used HMDB pathway information to flag Warburg-effect regions, and
\citet{wittmann2025} (S2IsoMEr / METASPACE) provide the methodologically central recent
tool, bootstrapping over isomer/isobar ambiguity (sampling one candidate per ion across
iterations and aggregating) and drawing pathway sets from RAMP-DB
(SMPDB/Reactome/KEGG/WikiPathways). Our isobar finding (Section~\ref{sec:realdata}) is a
direct empirical motivation for exactly this bootstrap: three of our $27$ columns are
ambiguous, so any pathway label assigned to a component's top ions must propagate that
uncertainty. A known IMS confounder to flag in any downstream labelling is matrix-adduct
formation, which co-localises with parent metabolites and can populate a spatial factor with
an adduct series rather than biology \citep{janda2021}. More broadly, \citet{alexandrov2020}
notes that enrichment-to-pathways, routine in other omics, is barely tooled for IMS, where
only a small fraction of acquired signals are annotated.

\paragraph{Independence without leaving the model.} The HSIC penalty
(Sections~\ref{sec:independence} and~\ref{sec:ica}) reconciles a second tension --- between
NMF's correlated parts and ICA's independent sources --- without abandoning non-negativity
or multiplicative coupling. Where prior MSI work either ran NMF \emph{or} ICA
\citep{siy2007}, we obtain ICA-style MI-sense independence \emph{as a tunable penalty on a
non-negative, multiplicatively-correct factorisation}, with a dimensionless weight that
interpolates between the two extremes and a kernel-alignment diagnostic to certify the
result.

\paragraph{Generality beyond IMS.} The same linear-space, multiplicative-error formulation
applies unchanged to spatial transcriptomics (Section~\ref{sec:spatial}): counts are kept
linear (a linear library-size correction, never a log) and KL/IS supplies the error model,
so a pathway's components remain genuine multiplicative outer products. The
spatial-transcriptomics result also exercises a capability latent in the IMS analysis but
not displayed there---decomposing \emph{one} pathway into sub-programmes that dominate
different tissue regions, read off the fraction-explained map $G_{p,k}$---suggesting the
method's natural use is less ``one pathway, one map'' than ``one pathway, a set of
region-specific programmes.'' That counts and metabolite intensities are handled by the
identical engine, differing only in the noise-model diagnostic's choice of loss, is itself
evidence that the coherence argument, not any modality-specific trick, is doing the work.

\paragraph{Novelty, stated plainly.} The factorisation itself is not novel --- IS/KL-NMF,
masked NMF, and kernel independence contrasts are standard. The defensible methodological
core is the \textbf{error-vs-coupling resolution}: modelling pathway co-regulation as
multiplicative and matching it to a multiplicative error model (linear-space IS/KL), which
the default log-then-PCA pipeline gets wrong; the independence penalty is a clean, optional
add-on rather than a separate algorithm. The likely publishable increment is the combination
--- coupling a multiplicatively-correct, saturation-aware factorisation to
\textbf{bootstrapped, isomer-aware pathway enrichment} in which the enrichment ranking
statistic \emph{is} a component's loading vector (ranking ions by loading on spatial factor
$k$). Existing enrichment tooling targets whole-dataset over-representation or single-cell
ranking, not ``rank ions by loading on a spatial component,'' and the intersection of
(multiplicative factorisation) $\times$ (ambiguity-aware enrichment) $\times$ (spatial IMS)
appears unoccupied. The contribution should be framed as the pipeline and its statistical
coherence, not any single algorithm.

\paragraph{Limitations.} (i)~The noise exponent here ($1.35$) is intermediate; KL is a
pragmatic choice and the IS/KL boundary is soft --- a per-ion or mixture noise model would be
more principled. (ii)~KL-NMF on raw large counts converges slowly; our reported fit had not
reached the tolerance at $800$ iterations, though the reconstruction was already stable and
the spatial structure is robust to further iteration. (iii)~$K$ selection by an elbow
heuristic is unstable in general; stability-based or cross-validated selection is preferable.
(iv)~Results are shown on one section and one synthetic dataset; spatial adjacency,
multi-sample alignment, adduct confounding, and the full enrichment step remain to be
evaluated. (v)~The independence penalty's random-Fourier-feature bandwidth is fixed per
iteration from the current factor, so its objective is only piecewise-deterministic across
iterations.

\paragraph{Conclusion.} Posing IMS pathway decomposition as the joint satisfaction of a
multiplicative coupling model and a multiplicative error model leads to a simple, coherent
pipeline --- linear-space IS/KL NMF with a censoring mask, plus an optional ICA-style
independence penalty --- that recovers known structure on synthetic data and produces
spatially distinct, chemically coherent components on real MALDI-MSI, while transparently
exposing the isomer ambiguity that any honest pathway labelling must confront.

\section*{Methods availability}

The decomposition engine, illustration layer, diagnostics, synthetic-data demo, and unit
tests are implemented in the \code{GAIT} package (\code{gait/decomposition/},
\code{gait/illustration/}). The end-to-end pipeline of Section~\ref{sec:results} is
reproduced by \code{demo\_decomposition.py}; design details are documented in
\code{CLAUDE.md}. The spatial-transcriptomics extension (Section~\ref{sec:spatial}) lives in
\code{gait/spatial.py} (10x Space~Ranger \code{.h5} loading, linear library-size
normalisation, gene-set selection, and the dominant-component view) and is reproduced by
\code{demo\_visium\_dopamine.py}.

\bibliography{ref}

\end{document}
